% DELTA23 agent /root/pentagon. New finite-quotient obstruction, not A.84 decision.
\subsection{Independent central flips in all four-point forgetting quotients}

\paragraph{D23.P1 (independent central coordinates).}
Let $P_n=\widehat\Gamma_{0,n}$, $n\geq4$. For each four-element
subset $I$ of the labelled punctures, let $\pi_I:P_n\to P_4$ be
forgetting the other punctures, with any orientation-preserving marking
of the target. Use the signed modular quotient of LM7\_PUSH E.58.11,
\[
 q:P_4\cong\widehat F_2\longrightarrow
 V_4\times\operatorname{SL}_2(\mathbb F_3),
\]
whose level-twelve projective notation sends $(1,-I)$ to $[5I]$.
Write $z=(1,-I)$. Then the simultaneous map has exact image
\[
 (q\pi_I)_I(P_n'')=\prod_{|I|=4}\langle z\rangle.
 \tag{D23.P1.1}
\]
Here the primes denote closed derived subgroups. In particular the
central signs in distinct four-point quotients can be prescribed
independently by one actual element of $P_n''$.

\begin{proof}
The second derived subgroup of $V_4\times\operatorname{SL}_2(\mathbb F_3)$
is $\{1\}\times\{I,-I\}$. The relevant explicit matrices are
\[
 x=\begin{pmatrix}1&2\\0&1\end{pmatrix},\qquad
 y=\begin{pmatrix}1&0\\1&1\end{pmatrix}.
\]
They generate $\operatorname{SL}_2(\mathbb F_3)$; its commutator
subgroup is the quaternion group of order eight and that group's
commutator subgroup is $\{I,-I\}$. These finite assertions, as well
as the matrix calculation below, are enumerated in the certificate
\path{agents/pentagon/check_four_point_central_flips.py}.
This gives the inclusion from left to right in (D23.P1.1).

Fix $I=\{i_1,i_2,i_3,i_4\}$. Choose a disk neighbourhood of a
chain joining these four punctures, with no other punctures in it.
Let $a,b,c\in P_n$ be the positive twists about the consecutive
pairs $(i_1,i_2),(i_2,i_3),(i_3,i_4)$ in that disk, and set
\[
 w_I=[[a,b],a[b,c]a^{-1}],\qquad [u,v]=uvu^{-1}v^{-1}.
 \tag{D23.P1.2}
\]
This lies in $P_n''$. If any puncture of $I$ is forgotten, one of
the two inner commutators is trivial: forgetting $i_1$ kills $a$,
forgetting $i_4$ kills $c$, and forgetting $i_2$ or $i_3$ kills
both twists in at least one inner commutator. A twist about a disk
containing at most one retained puncture is trivial in the punctured
sphere mapping class group. Therefore $\pi_J(w_I)=1$ for every
four-element subset $J\ne I$.

After forgetting all punctures outside $I$, the curves defining $a$
and $c$ are complementary and isotopic in the four-punctured sphere.
Thus $\pi_I(c)=\pi_I(a)$, and in the chosen standard marking their
signed modulo-three images are $q_3(a)=q_3(c)=x$, $q_3(b)=y$.
Direct multiplication gives
\[
 [x,y]=\begin{pmatrix}1&2\\2&2\end{pmatrix},\quad
 x[y,x]x^{-1}=\begin{pmatrix}1&1\\1&2\end{pmatrix},\quad
 [[x,y],x[y,x]x^{-1}]=-I.
\]
The modulo-four coordinate is trivial because $w_I\in P_n''$ and
the modulo-four target is abelian. Hence $q\pi_I(w_I)=z$.
Any other orientation-preserving marking is obtained by a mapping
class on the four-punctured sphere. Its conjugation action preserves
the modular quotient and its central element $z$; equivalently one
uses the quotient $\Gamma_{0,[4]}\to C_2\amalg C_3$ and its normal
level-twelve kernel. Thus the value $z$ does not depend on the
chosen chain marking. The images of the $w_I$ are the coordinate
basis of the product in (D23.P1.1). Their finite products give every
prescribed tuple. This proves the reverse inclusion.
\end{proof}

\paragraph{D23.P2 (finite simultaneous frames do not remove the cubic flip).}
Let a finite subgroup $H$ of $\widehat\Gamma_{0,[n]}$ act on the
four-element subsets by its label permutation, and retain the finite
quotient obtained from all the preceding four-point modular maps
and the label permutation map. For every $H$-orbit $\mathcal O$ of
four-element subsets there is an element of the image of $P_n''$
that centralizes the image of $H$ and whose four-point coordinates
are $z$ on $\mathcal O$ and $1$ elsewhere.
Consequently, whenever a finite simultaneous frame for $H$ is used to
read an equianharmonic frame sign at a subset $I\in\mathcal O$,
this centralizer changes that cubic sign by one. At any four-point
harmonic coordinate it leaves the harmonic parity unchanged. It also
leaves the abelianization of the simultaneous pure frame unchanged.

\begin{proof}
Conjugation by an element of the full mapping class group transports
forgetting maps according to its permutation of the labels. On the
second-derived central coordinates its effect is exactly that
permutation: $z$ is the distinguished nonidentity element of a group
of order two and is fixed by changes of target marking. By D23.P1,
the orbit-indicator tuple is the image of a product of the $w_I$.
It is fixed by conjugation by every element of $H$, hence centralizes
its image in the stated finite quotient. Multiplying any chosen
simultaneous frame by this element preserves all its finite
conjugator equations.

On the cubic coordinate, $z$ is $-I$ in
$C_{\operatorname{SL}_2(\mathbb F_3)}(R_0)\cong C_6$; its sign
on the rational fixed line is $-1$, so $a_3$ changes by one.
On the harmonic coordinate, $z=-I=S^2$ in
$\langle S\rangle\cong C_4$; the parity of the exponent of $S$
does not change. Trivial coordinates do not change that parity
either. Since the multiplying element lies in $P_n''$, the pure
frame's abelianization is also unchanged.
\end{proof}

\paragraph{Exact scope.}
D23.P2 is a statement about the specified finite quotient and finite
conjugator equations. It does not construct a different actual
unique frame in the full profinite group, nor a different gentle
or GT element. In particular, the additional affine five-point
quotient used in DELTA\_ON\_GT FP.32 records more than these
four-point projections and can exclude this ambiguity by its
relation $a_3=\kappa_d\bmod2$. No claim is made that D23.P2
survives that additional quotient. This computation therefore
blocks one simultaneous-frame shortcut but does not decide A.84.
